\documentclass[aspectratio=169,10pt]{beamer}
% ---------------------------------------------------------------
% R3.07 SQL dans un langage de programmation -- preambule commun des diapositives
% Charte derivee de R5.A.08 Qualite de developpement (metropolis, dark tech)
% Palette IUT partagee (violet du logo + accent or)
% Utilisation :
%   \documentclass[aspectratio=169,10pt]{beamer}
%   % ---------------------------------------------------------------
% R3.07 SQL dans un langage de programmation -- preambule commun des diapositives
% Charte derivee de R5.A.08 Qualite de developpement (metropolis, dark tech)
% Palette IUT partagee (violet du logo + accent or)
% Utilisation :
%   \documentclass[aspectratio=169,10pt]{beamer}
%   % ---------------------------------------------------------------
% R3.07 SQL dans un langage de programmation -- preambule commun des diapositives
% Charte derivee de R5.A.08 Qualite de developpement (metropolis, dark tech)
% Palette IUT partagee (violet du logo + accent or)
% Utilisation :
%   \documentclass[aspectratio=169,10pt]{beamer}
%   \input{../../latex/common/r307-slides-preamble.tex}
%   \graphicspath{{../../latex/common/assets/}{img/}{figs/}}
% ---------------------------------------------------------------
\usetheme{metropolis}
\usepackage{appendixnumberbeamer}

\usepackage[T1]{fontenc}
\usepackage[utf8]{inputenc}
\usepackage[french]{babel}
\usepackage[scale=0.94]{tgheros}
\renewcommand{\familydefault}{\sfdefault}
\usepackage{amsmath,amssymb,amsfonts}
\usepackage{graphicx}
\usepackage{booktabs}
\usepackage{multirow}
\usepackage{tikz}
\usetikzlibrary{shapes.geometric, shapes, arrows, arrows.meta, positioning, calc, fit}
\usepackage{xcolor}
\usepackage{listings}
\usepackage[skins,listings]{tcolorbox}
\usepackage{qrcode}
\usepackage{hyperref}

% ---------------------------------------------------------------
% PALETTE (dark-mode tech: near-black canvas, violet (logo IUT) + gold accents)
% ---------------------------------------------------------------
\definecolor{qabg}{HTML}{0E1016}      % canvas background
\definecolor{qapanel}{HTML}{181B24}   % card / block panel
\definecolor{qapanelalt}{HTML}{20242F}% secondary panel tint
\definecolor{qaline}{HTML}{2B303D}    % hairline rules / table rules
\definecolor{qatext}{HTML}{ECEEF2}    % primary text (off-white)
\definecolor{qamuted}{HTML}{8B93A6}   % secondary / caption text
\definecolor{qaaccent}{HTML}{D93ED7}  % primary accent -- violet du logo IUT, eclairci pour fond sombre
\definecolor{qaaccent2}{HTML}{F2B84B} % secondary accent -- or/ambre
\definecolor{qacodekw}{HTML}{C792EA}  % code : mots-cles
\definecolor{qacodestr}{HTML}{C3E88D} % code : chaines
\definecolor{qacodecom}{HTML}{7F8AA3} % code : commentaires
\definecolor{qacodeann}{HTML}{F2B84B} % code : annotations / javadoc

\setbeamercolor{background canvas}{bg=qabg}
\setbeamercolor{normal text}{fg=qatext}
\setbeamercolor{palette primary}{bg=qabg, fg=qamuted}
\setbeamercolor{frametitle}{bg=qabg, fg=qatext}
\setbeamercolor{progress bar}{fg=qaaccent, bg=qaline}
\setbeamercolor{alerted text}{fg=qaaccent2}
\setbeamercolor{block title}{bg=qapanel, fg=qaaccent}
\setbeamercolor{block body}{bg=qapanel, fg=qatext}
\setbeamercolor{block title alerted}{bg=qapanel, fg=qaaccent2}
\setbeamercolor{block body alerted}{bg=qapanel, fg=qatext}
\setbeamercolor{block title example}{bg=qapanelalt, fg=qaaccent}
\setbeamercolor{block body example}{bg=qapanelalt, fg=qatext}
\setbeamercolor{example text}{fg=qaaccent}
\setbeamercolor{itemize item}{fg=qaaccent}
\setbeamercolor{itemize subitem}{fg=qaaccent2}
\setbeamercolor{enumerate item}{fg=qaaccent}
\setbeamercolor{enumerate subitem}{fg=qaaccent2}
\setbeamercolor{caption}{fg=qamuted}
\setbeamercolor{page number in head/foot}{fg=qamuted, bg=qabg}
\setbeamercolor{title}{fg=qatext}
\setbeamercolor{subtitle}{fg=qamuted}
\setbeamercolor{author}{fg=qatext}
\setbeamercolor{institute}{fg=qamuted}
\setbeamercolor{date}{fg=qamuted}
\setbeamercolor{section in toc}{fg=qatext}
\setbeamercolor{subsection in toc}{fg=qamuted}

\setbeamertemplate{blocks}[rounded][shadow=false]
\setbeamertemplate{itemize item}{\raisebox{0.15em}{\scriptsize\textcolor{qaaccent}{\ensuremath{\blacktriangleright}}}}
\setbeamertemplate{itemize subitem}{\raisebox{0.1em}{\tiny\textcolor{qaaccent2}{\ensuremath{\bullet}}}}

\setbeamerfont{frametitle}{series=\bfseries, size=\Large}
\setbeamerfont{title}{series=\bfseries, size=\Large}
\setbeamerfont{subtitle}{size=\normalsize}
\setbeamerfont{institute}{size=\scriptsize}
\setbeamerfont{date}{size=\scriptsize}
\setbeamerfont{author}{size=\small}

\metroset{sectionpage=none, subsectionpage=none, numbering=fraction, progressbar=frametitle}
\setbeamertemplate{section in toc}[sections numbered]

% ---------------------------------------------------------------
% LISTINGS -- theme sombre (lstlisting reste utilisable tel quel)
% \lstcode : police + couleur de base, a reutiliser dans les
%            surcharges locales (ex. basicstyle=\lstcode\scriptsize)
% ---------------------------------------------------------------
\newcommand{\lstcode}{\ttfamily\color{qatext}}

\lstdefinelanguage{TypeScript}{
  keywords={typeof, new, true, false, catch, function, return, null, switch, var,
    if, in, while, do, else, case, break, const, let, interface, class, extends,
    implements, enum, export, import, this, super, constructor, public, private,
    protected, static, async, await, throw, throws, try, finally, void},
  ndkeywords={boolean, any, number, string, int, long, double},
  sensitive=true,
  comment=[l]{//},
  morecomment=[s]{/*}{*/},
  morestring=[b]',
  morestring=[b]"
}

\lstset{
  language=Java,
  basicstyle=\lstcode\footnotesize,
  keywordstyle=\color{qacodekw}\bfseries,
  ndkeywordstyle=\color{qaaccent2}\bfseries,
  identifierstyle=\color{qatext},
  stringstyle=\color{qacodestr},
  commentstyle=\color{qacodecom}\itshape,
  morecomment=[s][\color{qacodeann}]{/**}{*/},
  backgroundcolor=\color{qapanelalt},
  rulecolor=\color{qaline},
  frame=single,
  framesep=3pt,
  xleftmargin=4pt,
  xrightmargin=4pt,
  keepspaces=true,
  columns=fullflexible,
  showstringspaces=false,
  breaklines=true,
  tabsize=2,
  escapeinside={(*@}{@*)},
}

% ---------------------------------------------------------------
% CARD-STYLE BOXES
% ---------------------------------------------------------------
\tcbset{
  figcard/.style={colback=white, colframe=white, boxrule=0pt, arc=4pt,
    left=3pt, right=3pt, top=3pt, bottom=3pt, boxsep=0pt,
    halign=center, valign=center},
  datacard/.style={colback=qapanel, colframe=qaline, boxrule=0.4pt, arc=3pt,
    left=8pt, right=8pt, top=6pt, bottom=6pt, coltext=qatext,
    before upper={\color{qatext}}, halign=center},
  codebox/.style={colback=qapanelalt, colframe=qaline, boxrule=0.4pt, arc=3pt,
    left=8pt, right=8pt, top=5pt, bottom=5pt, listing only,
    listing options={basicstyle=\ttfamily\small\color{qaaccent},
      breaklines=true, showstringspaces=false, columns=fullflexible,
      keepspaces=true, frame=none, backgroundcolor={}}}
}
\newtcolorbox{figcard}[1][]{figcard,#1}
\newtcolorbox{datacard}{datacard}
\newtcblisting{codebox}{codebox}

% ---------------------------------------------------------------
% HANDY MACROS
% ---------------------------------------------------------------
\newcommand{\hl}[1]{\alert{#1}}
\newcommand{\source}[1]{{\scriptsize\color{qamuted}#1}}

% Sommaire automatique en debut de section
\AtBeginSection[]{%
  \begin{frame}<beamer>{Plan du cours}
    \tableofcontents[currentsection]
  \end{frame}}

% Page de titre : \rtitleslide{numero de chapitre}{titre du chapitre}
% Les auteurs sont repris de \rauthors (modifiable avant \begin{document})
\newcommand{\rauthors}{Selena Lamari et Chouki Tibermacine}
\newcommand{\rtitleslide}[2]{%
\begin{frame}[plain,noframenumbering]
  \begin{center}
    \vspace{2mm}
    {\color{qaaccent}\rule{2.1cm}{1.2pt}}\\[6mm]
    {\Large\bfseries\color{qatext} R3.07}\\[1mm]
    {\Large\bfseries\color{qatext} SQL dans un langage de programmation}\\[5mm]
    {\large\color{qamuted} Chapitre #1}\\[1mm]
    {\large\color{qaaccent} #2}\\[6mm]
    \includegraphics[height=2.1cm]{iut_logo_transparent.png}

    \vspace{4mm}

    {\small\color{qatext} \rauthors}\\[1mm]
    {\scriptsize\color{qamuted} BUT 2 Informatique --- IUT de Vannes --- Universit\'e Bretagne Sud}
  \end{center}
\end{frame}
}

% ---- Diapo de fin « Questions ? » (même modèle que R5.A.08) ----
\newcommand{\rcontact}{Chouki.Tibermacine@univ-ubs.fr}
\newcommand{\rquestionslide}{%
\begin{frame}{Questions ?}
  \begin{center}
    \vspace{2mm}
    \begin{tikzpicture}
      \node[regular polygon, regular polygon sides=6, draw=qaline, thick,
        minimum size=4.6cm, rotate=90] (hexin) {};
      \node[regular polygon, regular polygon sides=6, draw=qaaccent, dashed,
        thick, minimum size=6.2cm, rotate=90] (hexout) {};
      \node[font=\fontsize{92}{92}\selectfont\bfseries, color=qaaccent2] at (hexin) {?};
    \end{tikzpicture}\\[5mm]
    {\small\color{qamuted} \rcontact}
  \end{center}
\end{frame}}

%   \graphicspath{{../../latex/common/assets/}{img/}{figs/}}
% ---------------------------------------------------------------
\usetheme{metropolis}
\usepackage{appendixnumberbeamer}

\usepackage[T1]{fontenc}
\usepackage[utf8]{inputenc}
\usepackage[french]{babel}
\usepackage[scale=0.94]{tgheros}
\renewcommand{\familydefault}{\sfdefault}
\usepackage{amsmath,amssymb,amsfonts}
\usepackage{graphicx}
\usepackage{booktabs}
\usepackage{multirow}
\usepackage{tikz}
\usetikzlibrary{shapes.geometric, shapes, arrows, arrows.meta, positioning, calc, fit}
\usepackage{xcolor}
\usepackage{listings}
\usepackage[skins,listings]{tcolorbox}
\usepackage{qrcode}
\usepackage{hyperref}

% ---------------------------------------------------------------
% PALETTE (dark-mode tech: near-black canvas, violet (logo IUT) + gold accents)
% ---------------------------------------------------------------
\definecolor{qabg}{HTML}{0E1016}      % canvas background
\definecolor{qapanel}{HTML}{181B24}   % card / block panel
\definecolor{qapanelalt}{HTML}{20242F}% secondary panel tint
\definecolor{qaline}{HTML}{2B303D}    % hairline rules / table rules
\definecolor{qatext}{HTML}{ECEEF2}    % primary text (off-white)
\definecolor{qamuted}{HTML}{8B93A6}   % secondary / caption text
\definecolor{qaaccent}{HTML}{D93ED7}  % primary accent -- violet du logo IUT, eclairci pour fond sombre
\definecolor{qaaccent2}{HTML}{F2B84B} % secondary accent -- or/ambre
\definecolor{qacodekw}{HTML}{C792EA}  % code : mots-cles
\definecolor{qacodestr}{HTML}{C3E88D} % code : chaines
\definecolor{qacodecom}{HTML}{7F8AA3} % code : commentaires
\definecolor{qacodeann}{HTML}{F2B84B} % code : annotations / javadoc

\setbeamercolor{background canvas}{bg=qabg}
\setbeamercolor{normal text}{fg=qatext}
\setbeamercolor{palette primary}{bg=qabg, fg=qamuted}
\setbeamercolor{frametitle}{bg=qabg, fg=qatext}
\setbeamercolor{progress bar}{fg=qaaccent, bg=qaline}
\setbeamercolor{alerted text}{fg=qaaccent2}
\setbeamercolor{block title}{bg=qapanel, fg=qaaccent}
\setbeamercolor{block body}{bg=qapanel, fg=qatext}
\setbeamercolor{block title alerted}{bg=qapanel, fg=qaaccent2}
\setbeamercolor{block body alerted}{bg=qapanel, fg=qatext}
\setbeamercolor{block title example}{bg=qapanelalt, fg=qaaccent}
\setbeamercolor{block body example}{bg=qapanelalt, fg=qatext}
\setbeamercolor{example text}{fg=qaaccent}
\setbeamercolor{itemize item}{fg=qaaccent}
\setbeamercolor{itemize subitem}{fg=qaaccent2}
\setbeamercolor{enumerate item}{fg=qaaccent}
\setbeamercolor{enumerate subitem}{fg=qaaccent2}
\setbeamercolor{caption}{fg=qamuted}
\setbeamercolor{page number in head/foot}{fg=qamuted, bg=qabg}
\setbeamercolor{title}{fg=qatext}
\setbeamercolor{subtitle}{fg=qamuted}
\setbeamercolor{author}{fg=qatext}
\setbeamercolor{institute}{fg=qamuted}
\setbeamercolor{date}{fg=qamuted}
\setbeamercolor{section in toc}{fg=qatext}
\setbeamercolor{subsection in toc}{fg=qamuted}

\setbeamertemplate{blocks}[rounded][shadow=false]
\setbeamertemplate{itemize item}{\raisebox{0.15em}{\scriptsize\textcolor{qaaccent}{\ensuremath{\blacktriangleright}}}}
\setbeamertemplate{itemize subitem}{\raisebox{0.1em}{\tiny\textcolor{qaaccent2}{\ensuremath{\bullet}}}}

\setbeamerfont{frametitle}{series=\bfseries, size=\Large}
\setbeamerfont{title}{series=\bfseries, size=\Large}
\setbeamerfont{subtitle}{size=\normalsize}
\setbeamerfont{institute}{size=\scriptsize}
\setbeamerfont{date}{size=\scriptsize}
\setbeamerfont{author}{size=\small}

\metroset{sectionpage=none, subsectionpage=none, numbering=fraction, progressbar=frametitle}
\setbeamertemplate{section in toc}[sections numbered]

% ---------------------------------------------------------------
% LISTINGS -- theme sombre (lstlisting reste utilisable tel quel)
% \lstcode : police + couleur de base, a reutiliser dans les
%            surcharges locales (ex. basicstyle=\lstcode\scriptsize)
% ---------------------------------------------------------------
\newcommand{\lstcode}{\ttfamily\color{qatext}}

\lstdefinelanguage{TypeScript}{
  keywords={typeof, new, true, false, catch, function, return, null, switch, var,
    if, in, while, do, else, case, break, const, let, interface, class, extends,
    implements, enum, export, import, this, super, constructor, public, private,
    protected, static, async, await, throw, throws, try, finally, void},
  ndkeywords={boolean, any, number, string, int, long, double},
  sensitive=true,
  comment=[l]{//},
  morecomment=[s]{/*}{*/},
  morestring=[b]',
  morestring=[b]"
}

\lstset{
  language=Java,
  basicstyle=\lstcode\footnotesize,
  keywordstyle=\color{qacodekw}\bfseries,
  ndkeywordstyle=\color{qaaccent2}\bfseries,
  identifierstyle=\color{qatext},
  stringstyle=\color{qacodestr},
  commentstyle=\color{qacodecom}\itshape,
  morecomment=[s][\color{qacodeann}]{/**}{*/},
  backgroundcolor=\color{qapanelalt},
  rulecolor=\color{qaline},
  frame=single,
  framesep=3pt,
  xleftmargin=4pt,
  xrightmargin=4pt,
  keepspaces=true,
  columns=fullflexible,
  showstringspaces=false,
  breaklines=true,
  tabsize=2,
  escapeinside={(*@}{@*)},
}

% ---------------------------------------------------------------
% CARD-STYLE BOXES
% ---------------------------------------------------------------
\tcbset{
  figcard/.style={colback=white, colframe=white, boxrule=0pt, arc=4pt,
    left=3pt, right=3pt, top=3pt, bottom=3pt, boxsep=0pt,
    halign=center, valign=center},
  datacard/.style={colback=qapanel, colframe=qaline, boxrule=0.4pt, arc=3pt,
    left=8pt, right=8pt, top=6pt, bottom=6pt, coltext=qatext,
    before upper={\color{qatext}}, halign=center},
  codebox/.style={colback=qapanelalt, colframe=qaline, boxrule=0.4pt, arc=3pt,
    left=8pt, right=8pt, top=5pt, bottom=5pt, listing only,
    listing options={basicstyle=\ttfamily\small\color{qaaccent},
      breaklines=true, showstringspaces=false, columns=fullflexible,
      keepspaces=true, frame=none, backgroundcolor={}}}
}
\newtcolorbox{figcard}[1][]{figcard,#1}
\newtcolorbox{datacard}{datacard}
\newtcblisting{codebox}{codebox}

% ---------------------------------------------------------------
% HANDY MACROS
% ---------------------------------------------------------------
\newcommand{\hl}[1]{\alert{#1}}
\newcommand{\source}[1]{{\scriptsize\color{qamuted}#1}}

% Sommaire automatique en debut de section
\AtBeginSection[]{%
  \begin{frame}<beamer>{Plan du cours}
    \tableofcontents[currentsection]
  \end{frame}}

% Page de titre : \rtitleslide{numero de chapitre}{titre du chapitre}
% Les auteurs sont repris de \rauthors (modifiable avant \begin{document})
\newcommand{\rauthors}{Selena Lamari et Chouki Tibermacine}
\newcommand{\rtitleslide}[2]{%
\begin{frame}[plain,noframenumbering]
  \begin{center}
    \vspace{2mm}
    {\color{qaaccent}\rule{2.1cm}{1.2pt}}\\[6mm]
    {\Large\bfseries\color{qatext} R3.07}\\[1mm]
    {\Large\bfseries\color{qatext} SQL dans un langage de programmation}\\[5mm]
    {\large\color{qamuted} Chapitre #1}\\[1mm]
    {\large\color{qaaccent} #2}\\[6mm]
    \includegraphics[height=2.1cm]{iut_logo_transparent.png}

    \vspace{4mm}

    {\small\color{qatext} \rauthors}\\[1mm]
    {\scriptsize\color{qamuted} BUT 2 Informatique --- IUT de Vannes --- Universit\'e Bretagne Sud}
  \end{center}
\end{frame}
}

% ---- Diapo de fin « Questions ? » (même modèle que R5.A.08) ----
\newcommand{\rcontact}{Chouki.Tibermacine@univ-ubs.fr}
\newcommand{\rquestionslide}{%
\begin{frame}{Questions ?}
  \begin{center}
    \vspace{2mm}
    \begin{tikzpicture}
      \node[regular polygon, regular polygon sides=6, draw=qaline, thick,
        minimum size=4.6cm, rotate=90] (hexin) {};
      \node[regular polygon, regular polygon sides=6, draw=qaaccent, dashed,
        thick, minimum size=6.2cm, rotate=90] (hexout) {};
      \node[font=\fontsize{92}{92}\selectfont\bfseries, color=qaaccent2] at (hexin) {?};
    \end{tikzpicture}\\[5mm]
    {\small\color{qamuted} \rcontact}
  \end{center}
\end{frame}}

%   \graphicspath{{../../latex/common/assets/}{img/}{figs/}}
% ---------------------------------------------------------------
\usetheme{metropolis}
\usepackage{appendixnumberbeamer}

\usepackage[T1]{fontenc}
\usepackage[utf8]{inputenc}
\usepackage[french]{babel}
\usepackage[scale=0.94]{tgheros}
\renewcommand{\familydefault}{\sfdefault}
\usepackage{amsmath,amssymb,amsfonts}
\usepackage{graphicx}
\usepackage{booktabs}
\usepackage{multirow}
\usepackage{tikz}
\usetikzlibrary{shapes.geometric, shapes, arrows, arrows.meta, positioning, calc, fit}
\usepackage{xcolor}
\usepackage{listings}
\usepackage[skins,listings]{tcolorbox}
\usepackage{qrcode}
\usepackage{hyperref}

% ---------------------------------------------------------------
% PALETTE (dark-mode tech: near-black canvas, violet (logo IUT) + gold accents)
% ---------------------------------------------------------------
\definecolor{qabg}{HTML}{0E1016}      % canvas background
\definecolor{qapanel}{HTML}{181B24}   % card / block panel
\definecolor{qapanelalt}{HTML}{20242F}% secondary panel tint
\definecolor{qaline}{HTML}{2B303D}    % hairline rules / table rules
\definecolor{qatext}{HTML}{ECEEF2}    % primary text (off-white)
\definecolor{qamuted}{HTML}{8B93A6}   % secondary / caption text
\definecolor{qaaccent}{HTML}{D93ED7}  % primary accent -- violet du logo IUT, eclairci pour fond sombre
\definecolor{qaaccent2}{HTML}{F2B84B} % secondary accent -- or/ambre
\definecolor{qacodekw}{HTML}{C792EA}  % code : mots-cles
\definecolor{qacodestr}{HTML}{C3E88D} % code : chaines
\definecolor{qacodecom}{HTML}{7F8AA3} % code : commentaires
\definecolor{qacodeann}{HTML}{F2B84B} % code : annotations / javadoc

\setbeamercolor{background canvas}{bg=qabg}
\setbeamercolor{normal text}{fg=qatext}
\setbeamercolor{palette primary}{bg=qabg, fg=qamuted}
\setbeamercolor{frametitle}{bg=qabg, fg=qatext}
\setbeamercolor{progress bar}{fg=qaaccent, bg=qaline}
\setbeamercolor{alerted text}{fg=qaaccent2}
\setbeamercolor{block title}{bg=qapanel, fg=qaaccent}
\setbeamercolor{block body}{bg=qapanel, fg=qatext}
\setbeamercolor{block title alerted}{bg=qapanel, fg=qaaccent2}
\setbeamercolor{block body alerted}{bg=qapanel, fg=qatext}
\setbeamercolor{block title example}{bg=qapanelalt, fg=qaaccent}
\setbeamercolor{block body example}{bg=qapanelalt, fg=qatext}
\setbeamercolor{example text}{fg=qaaccent}
\setbeamercolor{itemize item}{fg=qaaccent}
\setbeamercolor{itemize subitem}{fg=qaaccent2}
\setbeamercolor{enumerate item}{fg=qaaccent}
\setbeamercolor{enumerate subitem}{fg=qaaccent2}
\setbeamercolor{caption}{fg=qamuted}
\setbeamercolor{page number in head/foot}{fg=qamuted, bg=qabg}
\setbeamercolor{title}{fg=qatext}
\setbeamercolor{subtitle}{fg=qamuted}
\setbeamercolor{author}{fg=qatext}
\setbeamercolor{institute}{fg=qamuted}
\setbeamercolor{date}{fg=qamuted}
\setbeamercolor{section in toc}{fg=qatext}
\setbeamercolor{subsection in toc}{fg=qamuted}

\setbeamertemplate{blocks}[rounded][shadow=false]
\setbeamertemplate{itemize item}{\raisebox{0.15em}{\scriptsize\textcolor{qaaccent}{\ensuremath{\blacktriangleright}}}}
\setbeamertemplate{itemize subitem}{\raisebox{0.1em}{\tiny\textcolor{qaaccent2}{\ensuremath{\bullet}}}}

\setbeamerfont{frametitle}{series=\bfseries, size=\Large}
\setbeamerfont{title}{series=\bfseries, size=\Large}
\setbeamerfont{subtitle}{size=\normalsize}
\setbeamerfont{institute}{size=\scriptsize}
\setbeamerfont{date}{size=\scriptsize}
\setbeamerfont{author}{size=\small}

\metroset{sectionpage=none, subsectionpage=none, numbering=fraction, progressbar=frametitle}
\setbeamertemplate{section in toc}[sections numbered]

% ---------------------------------------------------------------
% LISTINGS -- theme sombre (lstlisting reste utilisable tel quel)
% \lstcode : police + couleur de base, a reutiliser dans les
%            surcharges locales (ex. basicstyle=\lstcode\scriptsize)
% ---------------------------------------------------------------
\newcommand{\lstcode}{\ttfamily\color{qatext}}

\lstdefinelanguage{TypeScript}{
  keywords={typeof, new, true, false, catch, function, return, null, switch, var,
    if, in, while, do, else, case, break, const, let, interface, class, extends,
    implements, enum, export, import, this, super, constructor, public, private,
    protected, static, async, await, throw, throws, try, finally, void},
  ndkeywords={boolean, any, number, string, int, long, double},
  sensitive=true,
  comment=[l]{//},
  morecomment=[s]{/*}{*/},
  morestring=[b]',
  morestring=[b]"
}

\lstset{
  language=Java,
  basicstyle=\lstcode\footnotesize,
  keywordstyle=\color{qacodekw}\bfseries,
  ndkeywordstyle=\color{qaaccent2}\bfseries,
  identifierstyle=\color{qatext},
  stringstyle=\color{qacodestr},
  commentstyle=\color{qacodecom}\itshape,
  morecomment=[s][\color{qacodeann}]{/**}{*/},
  backgroundcolor=\color{qapanelalt},
  rulecolor=\color{qaline},
  frame=single,
  framesep=3pt,
  xleftmargin=4pt,
  xrightmargin=4pt,
  keepspaces=true,
  columns=fullflexible,
  showstringspaces=false,
  breaklines=true,
  tabsize=2,
  escapeinside={(*@}{@*)},
}

% ---------------------------------------------------------------
% CARD-STYLE BOXES
% ---------------------------------------------------------------
\tcbset{
  figcard/.style={colback=white, colframe=white, boxrule=0pt, arc=4pt,
    left=3pt, right=3pt, top=3pt, bottom=3pt, boxsep=0pt,
    halign=center, valign=center},
  datacard/.style={colback=qapanel, colframe=qaline, boxrule=0.4pt, arc=3pt,
    left=8pt, right=8pt, top=6pt, bottom=6pt, coltext=qatext,
    before upper={\color{qatext}}, halign=center},
  codebox/.style={colback=qapanelalt, colframe=qaline, boxrule=0.4pt, arc=3pt,
    left=8pt, right=8pt, top=5pt, bottom=5pt, listing only,
    listing options={basicstyle=\ttfamily\small\color{qaaccent},
      breaklines=true, showstringspaces=false, columns=fullflexible,
      keepspaces=true, frame=none, backgroundcolor={}}}
}
\newtcolorbox{figcard}[1][]{figcard,#1}
\newtcolorbox{datacard}{datacard}
\newtcblisting{codebox}{codebox}

% ---------------------------------------------------------------
% HANDY MACROS
% ---------------------------------------------------------------
\newcommand{\hl}[1]{\alert{#1}}
\newcommand{\source}[1]{{\scriptsize\color{qamuted}#1}}

% Sommaire automatique en debut de section
\AtBeginSection[]{%
  \begin{frame}<beamer>{Plan du cours}
    \tableofcontents[currentsection]
  \end{frame}}

% Page de titre : \rtitleslide{numero de chapitre}{titre du chapitre}
% Les auteurs sont repris de \rauthors (modifiable avant \begin{document})
\newcommand{\rauthors}{Selena Lamari et Chouki Tibermacine}
\newcommand{\rtitleslide}[2]{%
\begin{frame}[plain,noframenumbering]
  \begin{center}
    \vspace{2mm}
    {\color{qaaccent}\rule{2.1cm}{1.2pt}}\\[6mm]
    {\Large\bfseries\color{qatext} R3.07}\\[1mm]
    {\Large\bfseries\color{qatext} SQL dans un langage de programmation}\\[5mm]
    {\large\color{qamuted} Chapitre #1}\\[1mm]
    {\large\color{qaaccent} #2}\\[6mm]
    \includegraphics[height=2.1cm]{iut_logo_transparent.png}

    \vspace{4mm}

    {\small\color{qatext} \rauthors}\\[1mm]
    {\scriptsize\color{qamuted} BUT 2 Informatique --- IUT de Vannes --- Universit\'e Bretagne Sud}
  \end{center}
\end{frame}
}

% ---- Diapo de fin « Questions ? » (même modèle que R5.A.08) ----
\newcommand{\rcontact}{Chouki.Tibermacine@univ-ubs.fr}
\newcommand{\rquestionslide}{%
\begin{frame}{Questions ?}
  \begin{center}
    \vspace{2mm}
    \begin{tikzpicture}
      \node[regular polygon, regular polygon sides=6, draw=qaline, thick,
        minimum size=4.6cm, rotate=90] (hexin) {};
      \node[regular polygon, regular polygon sides=6, draw=qaaccent, dashed,
        thick, minimum size=6.2cm, rotate=90] (hexout) {};
      \node[font=\fontsize{92}{92}\selectfont\bfseries, color=qaaccent2] at (hexin) {?};
    \end{tikzpicture}\\[5mm]
    {\small\color{qamuted} \rcontact}
  \end{center}
\end{frame}}

\graphicspath{{../../latex/common/assets/}}
\usepackage{dirtree}

\title{Patrons DAO et DTO}
\author[C. Tibermacine]{Selena Lamari et Chouki Tibermacine}
\date{}

\begin{document}
	
	\rtitleslide{5}{Patrons DAO et DTO}


\begin{frame}{Introduction}
    \begin{block}{DAO (Data Access Object)}
        \begin{itemize}
            \item Séparation de la logique métier et de l'accès aux données.
            \item Fournit une interface unifiée pour accéder à une source de données (BDD, fichiers, etc.).
        \end{itemize}
    \end{block}
    \begin{block}{DTO (Data Transfer Object)}
        \begin{itemize}
            \item Objet simple utilisé pour transférer des données entre processus ou couches.
            \item Réduit le nombre d'appels réseau en regroupant les données.
        \end{itemize}
    \end{block}
\end{frame}

\begin{frame}[fragile]{Exemple DAO en Java (JDBC)}
    \begin{exampleblock}{Interface DAO}
        \begin{lstlisting}[language=Java]
public interface UserDAO {
    void addUser(User user);
    User getUser(int id);
    List<User> getAllUsers();
}
        \end{lstlisting}
    \end{exampleblock}
\end{frame}

\begin{frame}[fragile]{Exemple DAO en Java (Implémentation JDBC)}
        \begin{lstlisting}[language=Java]
public class UserDAOImpl implements UserDAO {
 private Connection connection;
 public UserDAOImpl(Connection connection) {
  this.connection = connection;
 }
 @Override
 public void addUser(User user) throws SQLException {
  String sql = "INSERT INTO users (name, email) VALUES (?, ?)";
  try (PreparedStatement stmt = connection.prepareStatement(sql)) {
   stmt.setString(1, user.getName());
   stmt.setString(2, user.getEmail());
   stmt.executeUpdate();
  }
 }
}
\end{lstlisting}
\end{frame}

\begin{frame}[fragile]{Exemple DTO en Java}
        \begin{lstlisting}[language=Java]
public class UserDTO {
    private int id;
    private String name;
    private String email;

    // Getters et Setters
    public int getId() { return id; }
    public void setId(int id) { this.id = id; }
    public String getName() { return name; }
    public void setName(String name) { this.name = name; }
    public String getEmail() { return email; }
    public void setEmail(String email) { this.email = email; }
}
        \end{lstlisting}
\end{frame}


\begin{frame}{Où utilise-t-on DAO et DTO ?}
    \begin{block}{DAO (Data Access Object)}
        \begin{itemize}
            \item Utilisé dans la \textbf{couche d'accès aux données}.
            \item Abstraction de la source de données (BDD, API, fichiers, etc.).
            \item Exemple : Accès à une base de données avec JDBC, Hibernate, ou un ORM.
        \end{itemize}
    \end{block}
\end{frame}

\begin{frame}{Où utilise-t-on DAO et DTO ? -suite-}    
    \begin{block}{DTO (Data Transfer Object)}
        \begin{itemize}
            \item Utilisé pour \textbf{transférer des données} entre couches (ex : entre la couche service et la couche présentation).
            \item Évite d'exposer les objets métier ou les entités issues de la bdd directement.
            \item Exemple : retourner un objet User avec des informations non-sensibles (retirer le password, ...).
            \item Exemple : réponse d'une API REST (alléger les réponses), données envoyées à une vue.
        \end{itemize}
    \end{block}
\end{frame}


\begin{frame}[fragile]{O\`u placer la fronti\`ere de transaction ?}
  \begin{itemize}
    \item Le code JDBC des transactions (\texttt{setAutoCommit}, \texttt{commit}, \texttt{rollback}) a \'et\'e vu au \alert{chapitre 2}
    \item Un DAO = une table/entit\'e. Or un cas d'usage m\'etier touche souvent \alert{plusieurs DAO}
    \item[$\Rightarrow$] La transaction se d\'emarre et se termine dans la \textbf{couche service}, et les DAO partagent la m\^eme connexion
  \end{itemize}
\begin{lstlisting}[language=Java]
public void virer(int src, int dst, double montant) throws SQLException {
  try (Connection c = dataSource.getConnection()) {
    c.setAutoCommit(false);
    try {
      compteDAO.debiter(c, src, montant);
      compteDAO.crediter(c, dst, montant);
      operationDAO.journaliser(c, src, dst, montant);
      c.commit();
    } catch (SQLException e) { c.rollback(); throw e; }
  }
}
\end{lstlisting}
  {\small\color{qamuted} Avec JPA ou Spring : m\^eme principe (\texttt{EntityTransaction}, puis \texttt{@Transactional} sur le service).}
\end{frame}

\begin{frame}[fragile]{Gestion des exceptions}
\textbf{Ne pas exposer} les exceptions SQL dans les couches supérieures.
Utiliser des exceptions métiers ou des DTOs pour retourner des erreurs
\begin{lstlisting}[language=Java]  
public UserDTO getUserById(int id) {
    try {
        // Logique SQL
    } catch (SQLException e) {
        throw new DataAccessException("Erreur lors de la recuperation de l'utilisateur", e);
    }
}

\end{lstlisting}
\end{frame}


\begin{frame}{Structure typique d'une application}
\begin{center}
    \begin{tikzpicture}[
        layer/.style={draw=qaaccent, rounded corners, fill=qapanel, text=qatext, font=\small, rectangle, minimum width=3cm, minimum height=0.7cm, align=center},
        node distance=0.6cm,
        arrow/.style={->, >=stealth', thick, qamuted}
    ]
        \node[layer, fill=qapanel] (presentation) {Présentation (API/UIView)};
        \node[layer, fill=qapanel, below=of presentation] (service) {Service (Logique métier)};
        \node[layer, fill=qapanel, below=of service] (dao) {DAO (Accès aux données)};
        \node[layer, fill=qapanel, below=of dao] (database) {Base de données};

        \draw[arrow] (presentation) -- node[anchor=west] {DTO} (service);
        \draw[arrow] (service) -- node[anchor=west] {DTO/Entité} (dao);
        \draw[arrow] (dao) -- node[anchor=west] {SQL/Requêtes} (database);
        \draw[arrow] (dao) -- node[anchor=east] {Entité} (service);
    \end{tikzpicture}
\end{center}
    \begin{center}
        \small{Les DTO circulent entre les couches, les DAO interagissent avec la bdd.}
    \end{center}
\end{frame}

\begin{frame}{Exemple d'architecture avec DAO et DTO}
\begin{center}
    \begin{tikzpicture}[
        layer/.style={draw=qaaccent, rounded corners, fill=qapanel, text=qatext, font=\small, rectangle, minimum width=2.5cm, minimum height=0.65cm, align=center},
        node distance=0.75cm and 1.8cm,
        arrow/.style={->, >=stealth', thick, qamuted}
    ]
        \node[layer, fill=qapanel] (api) {API REST};
        \node[layer, fill=qapanel, below=of api] (service) {Service};
        \node[layer, fill=qapanel, below=of service] (dao) {DAO};
        \node[layer, fill=qapanel, below=of dao] (db) {Base de données};

        \node[layer, right=of service, draw=qaaccent2, fill=qapanelalt] (dto) {DTO};
        \node[layer, right=of dao, draw=qaaccent2, fill=qapanelalt] (entity) {Entité JPA};

        \draw[arrow] (api) -- node[above] {requête} (service);
        \draw[arrow] (service) -- node[above] {DTO} (dto);
        \draw[arrow] (service) -- node[left] {appel} (dao);
        \draw[arrow] (dao) -- node[left] {requête} (db);
        \draw[arrow] (db) -- node[right] {résultat} (dao);
        \draw[arrow] (dao) -- node[above] {Entité} (entity);
        \draw[arrow] (entity) -- node[right] {mapping} (dto);
        \draw[arrow] (dto) -- node[above] {réponse} (api);
    \end{tikzpicture}
    \end{center}
    \begin{center}
        \small{Les DTO sont utilisés pour les échanges entre l'API et le service, et entre le service et la vue.}
    \end{center}
\end{frame}

\begin{frame}{Cas d'usage concret : Application Web}
    \begin{block}{Scénario}
        \begin{itemize}
            \item Un utilisateur demande la liste des produits via une API REST.
            \item Le \textbf{Service} appelle le \textbf{DAO} pour récupérer les données.
            \item Le \textbf{DAO} interroge la base de données et retourne une \textbf{Entité}.
            \item Le \textbf{Service} convertit l'Entité en \textbf{DTO} et renvoie la réponse.
        \end{itemize}
    \end{block}
\end{frame}

\begin{frame}{Cas d'usage concret : Application Web}    
    \begin{tikzpicture}[
        layer/.style={draw=qaaccent, rounded corners, fill=qapanel, text=qatext, font=\small, rectangle, minimum width=2.5cm, minimum height=0.8cm, align=center},
        arrow/.style={->, >=stealth', thick, qamuted}
    ]
        % Première ligne : Client et API
        \node[layer, fill=qapanel] (client) {Client (Navigateur)};
        \node[layer, fill=qapanel, right=3cm of client] (api) {API};

        % Deuxième ligne : Service, DAO, BDD
        \node[layer, fill=qapanel, below=1.5cm of api] (service) {Service};
        \node[layer, fill=qapanel, left=3.5cm of service] (dao) {DAO};
        \node[layer, fill=qapanel, below=2cm of dao] (db) {BDD};

        % Flèches entre Client et API (première ligne)
        \draw[arrow] (client) -- node[above] {GET /products} (api);

        % Flèches de API vers Service (verticale)
        \draw[arrow] (api) -- node[right] {getAllProducts()} (service);

        % Flèches entre Service, DAO et BDD (deuxième ligne)
        \draw[arrow] (service) -- node[above] {findAll()} (dao);
        \draw[arrow] (dao) -- node[right] {SQL} (db);

        % Flèches de retour
        \draw[arrow] (db) -- node[left] {Produits (Entité)} (dao);
        \draw[arrow] (dao) -- node[below] {Produits (Entité)} (service);
        \draw[arrow] (service) -- node[left] {Produits (DTO)} (api);
        \draw[arrow] (api) -- node[below] {JSON (DTO)} (client);
    \end{tikzpicture}
\end{frame}


\begin{frame}{Avantages de cette structure}
    \begin{columns}
        \column{0.5\textwidth}
        \begin{block}{DAO}
            \begin{itemize}
                \item Découplage de la logique métier et de l'accès aux données.
                \item Facilite les tests unitaires (mock des DAO).
                \item Permet de changer de source de données sans impacter le reste.
            \end{itemize}
        \end{block}
        \column{0.5\textwidth}
        \begin{block}{DTO}
            \begin{itemize}
                \item Contrôle précis des données exposées.
                \item Réduction du trafic réseau (seulement les données nécessaires).
                \item Sécurité (masquage des données sensibles).
            \end{itemize}
        \end{block}
    \end{columns}
    \begin{center}
        \small{DAO + DTO = Architecture modulaire, maintenable et sécurisée.}
    \end{center}
\end{frame}

\begin{frame}[fragile]{Structure de dossiers pour une application Java}
  \begin{columns}[T]
    \begin{column}{0.42\textwidth}
\begin{lstlisting}[language={},basicstyle=\lstcode\scriptsize]
src/
|-- main/
|   |-- java/com/monapplication/
|   |   |-- config/
|   |   |-- controller/
|   |   |-- service/
|   |   |-- model/
|   |   |   |-- entity/
|   |   |   |-- dao/
|   |   |   |-- dto/
|   |   |   `-- exception/
|   |   |-- util/
|   |   `-- Application.java
|   `-- resources/
|       |-- application.properties
|       `-- db/
`-- test/
\end{lstlisting}
    \end{column}
    \begin{column}{0.54\textwidth}
      \begin{exampleblock}{Explications}
        \scriptsize
\begin{itemize}
                    \item \texttt{config/}: Configuration de l'application.
                    \item \texttt{controller/}: Contrôleurs pour les requêtes HTTP.
                    \item \texttt{service/}: Logique métier.
                    \item \texttt{model/entity/}: Entités JPA.
                    \item \texttt{model/dao/}: Accès aux données.
                    \item \texttt{model/dto/}: Objets de transfert de données.
                    \item \texttt{model/exception/}: Exceptions personnalisées.
                    \item \texttt{util/}: Utilitaires et helpers.
                    \item \texttt{resources/}: Fichiers de configuration et scripts SQL.
                    \item \texttt{test/}: Tests unitaires et d'intégration.
                \end{itemize}
      \end{exampleblock}
    \end{column}
  \end{columns}
\end{frame}


\begin{frame}{Bonnes pratiques pour une architecture Java}
    \begin{block}{Séparation des couches}
        \begin{itemize}
            \item \textbf{Contrôleurs} : Gèrent les requêtes HTTP (API REST ou MVC).
            \item \textbf{Services} : Contiennent la \textbf{logique métier} et utilisent les DAO.
            \item \textbf{DAO} : Encapsulent les \textbf{requêtes SQL/JPA} et l'accès aux données.
            \item \textbf{DTO} : Transfèrent les données entre couches \textbf{sans logique métier}.
            \item \textbf{Entités} : Représentent les \textbf{tables de la base de données} (JPA).
        \end{itemize}
    \end{block}
\end{frame}


\begin{frame}{Bonnes pratiques pour une architecture Java -suite-}

    \begin{alertblock}{Règles à respecter}
        \begin{itemize}
            \item Ne \textbf{jamais} exposer les entités JPA directement dans une API.
            \item Utiliser des \textbf{DTO} pour les réponses API.
            \item Gérer les \textbf{transactions} au niveau des services.
            \item \textbf{Injecter les dépendances} (ex : \texttt{@Autowired} avec Spring).
            \item \textbf{Ne pas} appeler un DAO directement depuis un contrôleur.
        \end{itemize}
    \end{alertblock}

    \begin{exampleblock}{Exemple de flux}
        \begin{tikzpicture}[
            layer/.style={draw=qaaccent, rounded corners, fill=qapanel, text=qatext, font=\small, rectangle, minimum width=2cm, minimum height=0.6cm, align=center},
            arrow/.style={->, >=stealth', thick, qamuted}
        ]
            \node[layer, fill=qapanel] (controller) {Contrôleur};
            \node[layer, fill=qapanel, right=of controller] (service) {Service};
            \node[layer, fill=qapanel, right=of service] (dao) {DAO};
            \node[layer, fill=qapanel, right=of dao] (db) {Bdd};
            \node[layer, below=of service, draw=qaaccent2, fill=qapanelalt] (dto) {DTO};

            \draw[arrow] (controller) -- node[above] {requête} (service);
            \draw[arrow] (service) -- node[above] {appel} (dao);
            \draw[arrow] (dao) -- node[above] {SQL} (db);
            \draw[arrow] (db) -- node[below] {Entité} (dao);
            \draw[arrow] (dao) -- node[below] {Entité} (service);
            \draw[arrow] (service) -- node[right] {DTO} (dto);
            \draw[arrow] (dto) -- node[below] {réponse} (controller);
        \end{tikzpicture}
    \end{exampleblock}
\end{frame}


\begin{frame}{Conclusion}
    \begin{itemize}
        \item DAO et DTO sont complémentaires pour une architecture propre.
        \item DAO gère l'accès aux données, DTO transporte les données.
        \item Utilisables dans de nombreux langages (Java, TypeScript, Python, etc.).
    \end{itemize}
\end{frame}

\begin{frame}{Questions pour vous, alors}
\begin{enumerate}
\item Pourquoi utiliser un DAO plutôt que d’écrire des requêtes SQL directement dans le code métier ?
\item Comment un DTO peut-il améliorer les performances d’une application ?
\item Quels sont les risques de ne pas utiliser de transactions dans un DAO ?
\item Comment adapteriez-vous un DAO pour utiliser une base de données NoSQL (ex : MongoDB) ?

\end{enumerate}
\end{frame}

\rquestionslide

\end{document}
